\chapter{Stopping a ranked screen early}
\label{ch:stopping}

% Author's note (verification and reproduction): session screenlit-q1
% (and screenlit-q1-top); queries q1a-q1g, q1x, q1y; check by meaning
% checks/q1*.json; sample tier2-sample-q1.tsv (seed 20261124); see the
% block at the top of search.tex.  Sources: doc/bibliography.bib, keys
% Repke2026 Callaghan2020 Kempny2026 Campos2024 Koenig2024 Boetje2024
% Hamel2020 Gates2020 Callaghan2024 Slade2026 Briscoe2023.

The procedure of \cref{ch:procedure} screens a first tier, samples the
rest and states both; it never stops because the records have grown
thin.
Tools that rank records for screening offer another way: screen in the
order of the ranking and stop when a rule says that enough of the
relevant records have been found.
If such a rule were reliable, the first tier could end where the rule
says, and nothing more would need to be said about the rest.
This chapter asks whether it is: \emph{does stopping a prioritised
screen before the end of the list keep the share of missed relevant
studies small?}
The claim is tested at the strength a procedure would need: a rule that
holds a stated recall target across reviews, not one that sometimes
does.

\section{Method}
\label{sec:stopping-method}

The search is the one described in \cref{ch:lit-search}.
The first question's queries are listed below; \cref{tab:stopping-matrix}
gives their hits per database.
The support queries are q1a to q1e, the counter queries q1x and q1y; q1f
and q1g were added after the check by meaning, with the words that tool
evaluations use (sensitivity, missed records, workload; prioritisation,
filtering, ranking, work saved).
\begin{description}[style=nextline,font=\normalfont\bfseries]
\item[q1a, support query] {\footnotesize\ttfamily\raggedright (\textquotedbl{}stopping rule\textquotedbl{} OR \textquotedbl{}stopping rules\textquotedbl{} OR \textquotedbl{}stopping criterion\textquotedbl{} OR \textquotedbl{}stopping criteria\textquotedbl{} OR \textquotedbl{}stopping method\textquotedbl{}) AND (screening) AND (\textquotedbl{}systematic review\textquotedbl{} OR \textquotedbl{}technology-assisted review\textquotedbl{} OR \textquotedbl{}citation screening\textquotedbl{})\par}
\item[q1b, support query] {\footnotesize\ttfamily\raggedright (\textquotedbl{}when to stop\textquotedbl{} OR \textquotedbl{}stop screening\textquotedbl{} OR \textquotedbl{}early stopping\textquotedbl{}) AND (screening) AND (\textquotedbl{}systematic review\textquotedbl{} OR \textquotedbl{}technology-assisted review\textquotedbl{})\par}
\item[q1c, support query] {\footnotesize\ttfamily\raggedright (\textquotedbl{}active learning\textquotedbl{} OR \textquotedbl{}screening prioritisation\textquotedbl{} OR \textquotedbl{}screening prioritization\textquotedbl{} OR \textquotedbl{}prioritised screening\textquotedbl{}) AND (\textquotedbl{}systematic review\textquotedbl{} OR \textquotedbl{}citation screening\textquotedbl{}) AND (recall OR \textquotedbl{}work saved\textquotedbl{})\par}
\item[q1d, support query] {\footnotesize\ttfamily\raggedright (\textquotedbl{}capture-recapture\textquotedbl{} OR \textquotedbl{}work saved over sampling\textquotedbl{} OR \textquotedbl{}recall estimation\textquotedbl{}) AND (screening OR stopping) AND (\textquotedbl{}systematic review\textquotedbl{} OR \textquotedbl{}technology-assisted review\textquotedbl{})\par}
\item[q1e, support query] {\footnotesize\ttfamily\raggedright (\textquotedbl{}technology assisted review\textquotedbl{} OR \textquotedbl{}technology-assisted review\textquotedbl{} OR \textquotedbl{}machine learning\textquotedbl{} OR \textquotedbl{}active learning\textquotedbl{} OR \textquotedbl{}prioritised screening\textquotedbl{} OR \textquotedbl{}prioritized screening\textquotedbl{} OR \textquotedbl{}screening prioritisation\textquotedbl{} OR \textquotedbl{}screening prioritization\textquotedbl{}) AND (stopping OR stop OR \textquotedbl{}target recall\textquotedbl{} OR \textquotedbl{}recall target\textquotedbl{} OR \textquotedbl{}high recall\textquotedbl{} OR \textquotedbl{}total recall\textquotedbl{}) AND (\textquotedbl{}systematic review\textquotedbl{} OR \textquotedbl{}systematic reviews\textquotedbl{} OR \textquotedbl{}citation screening\textquotedbl{} OR \textquotedbl{}abstract screening\textquotedbl{})\par}
\item[q1f, support query] {\footnotesize\ttfamily\raggedright (\textquotedbl{}machine learning\textquotedbl{} OR \textquotedbl{}active learning\textquotedbl{} OR \textquotedbl{}semi-automated\textquotedbl{} OR \textquotedbl{}semi-automation\textquotedbl{} OR \textquotedbl{}automated screening\textquotedbl{} OR \textquotedbl{}screening prioritization\textquotedbl{} OR \textquotedbl{}screening prioritisation\textquotedbl{} OR \textquotedbl{}relevance prediction\textquotedbl{} OR \textquotedbl{}relevance predictions\textquotedbl{}) AND (\textquotedbl{}title and abstract screening\textquotedbl{} OR \textquotedbl{}abstract screening\textquotedbl{} OR \textquotedbl{}citation screening\textquotedbl{} OR \textquotedbl{}title screening\textquotedbl{}) AND (\textquotedbl{}systematic review\textquotedbl{} OR \textquotedbl{}systematic reviews\textquotedbl{} OR \textquotedbl{}rapid review\textquotedbl{}) AND (sensitivity OR recall OR missed OR \textquotedbl{}relevant studies\textquotedbl{} OR workload)\par}
\item[q1g, support query] {\footnotesize\ttfamily\raggedright (prioritization OR prioritisation OR prioritized OR prioritised OR filtering OR ranking) AND (\textquotedbl{}study selection\textquotedbl{} OR \textquotedbl{}title and abstract screening\textquotedbl{} OR \textquotedbl{}abstract screening\textquotedbl{} OR \textquotedbl{}citation screening\textquotedbl{}) AND (\textquotedbl{}systematic review\textquotedbl{} OR \textquotedbl{}systematic reviews\textquotedbl{}) AND (\textquotedbl{}work savings\textquotedbl{} OR \textquotedbl{}work saved\textquotedbl{} OR recall OR sensitivity OR missed OR stopping)\par}
\item[q1x, counter query] {\footnotesize\ttfamily\raggedright (\textquotedbl{}stopping rule\textquotedbl{} OR \textquotedbl{}stopping rules\textquotedbl{} OR \textquotedbl{}stopping criteria\textquotedbl{} OR \textquotedbl{}stopping method\textquotedbl{}) AND (unreliable OR fail OR missed OR premature OR undershoot OR unstable) AND (screening) AND (review)\par}
\item[q1y, counter query] {\footnotesize\ttfamily\raggedright (\textquotedbl{}stopping rule\textquotedbl{} OR \textquotedbl{}stopping rules\textquotedbl{} OR \textquotedbl{}stopping criteria\textquotedbl{} OR \textquotedbl{}stopping criterion\textquotedbl{} OR \textquotedbl{}stopping method\textquotedbl{} OR \textquotedbl{}when to stop\textquotedbl{} OR \textquotedbl{}early stopping\textquotedbl{} OR \textquotedbl{}active learning\textquotedbl{} OR \textquotedbl{}technology-assisted review\textquotedbl{} OR \textquotedbl{}screening prioritisation\textquotedbl{} OR \textquotedbl{}screening prioritization\textquotedbl{}) AND (missed OR miss OR fail OR failed OR unreliable OR overestimate OR \textquotedbl{}below target\textquotedbl{} OR \textquotedbl{}not reached\textquotedbl{} OR variability) AND (recall) AND (\textquotedbl{}systematic review\textquotedbl{} OR \textquotedbl{}systematic reviews\textquotedbl{} OR \textquotedbl{}citation screening\textquotedbl{})\par}
\end{description}


\begin{table}
  \begin{sidecaption}{%
      Hits per query and database for the first question (stopping a
      ranked screen early), searched 2026-11-23 with Scholar, as
      \texttt{scholar sessions show} reports them; the queries are
      listed in \cref{sec:stopping-method}.
      OpenAlex (OA; not open access) in title and abstract, Semantic
      Scholar (S2) through its bulk search, Scopus as
      \texttt{TITLE-ABS-KEY(\dots)}, Web of Science (WoS) as
      \texttt{TS=(\dots)}, IEEE Xplore (IEEE), arXiv; no limit, so the
      counts are all the records each database returned.
      \enquote{113 of about 142}: the database reported more records
      than it returned (Semantic Scholar's totals are approximate).
      DBLP was not searched (\cref{sec:search-databases}).
    }[tab:stopping-matrix]
    \footnotesize
    \begin{tabular}{@{}l*{6}{>{\raggedleft\arraybackslash}p{1.45cm}}@{}}
\toprule
Query & OA & S2 & Scopus & WoS & IEEE & arXiv \\
\midrule
q1a & 71 & 38 & 35 & 25 & 2 & 6 \\
q1b & 35 & 33 & 20 & 15 & 1 & 3 \\
q1c & 58 & 45 & 47 & 26 & 6 & 4 \\
q1d & 36 & 34 & 26 & 20 & 3 & 1 \\
q1e & 193 & 105 & 116 & 73 & 37 & 11 \\
q1f & 148 & 113 of about 142 & 123 & 86 & 8 & 7 \\
q1g & 224 & 127 & 91 & 117 & 5 & 3 \\
q1x & 29 & 10 & 7 & 9 & 0 & 0 \\
q1y & 14 & 8 & 4 & 6 & 0 & 0 \\
\bottomrule
\end{tabular}

  \end{sidecaption}
\end{table}

Five studies were named before the search: two evaluations of stopping
rules (one in education, one in psychology), the statistical stopping
criterion, an analysis of the work-saved measure, and the
capture-recapture stopping rule; all five came back.
The check by meaning returned 29 works that no session held; the two
repair queries brought back 17 of them (\cref{sec:search-check}).
Of the 931 distinct records, 454 were screened in the first tier and 53
of them bore on the question (23 supporting the claim, 30 qualifying or
contradicting it); of the 477 left unscreened, a random sample of 60
held 3 bearing records (5\,\%; 95\,\% interval 1.4 to 13.7\,\%).
All three are evaluations of screening tools, published in journals of
research-synthesis methods, read only in their abstracts and not cited;
they agree with the sources below.
The records judged bearing are listed in \cref{sec:stopping-hits}.

\section{Results}
\label{sec:stopping-results}

\paragraph{Most stopping rules miss their target.}
The largest comparison found tested 15 stopping methods on 81 datasets
of real reviews and found that
\textcquote[abstract]{Repke2026}{almost all existing stopping methods
either fail to reliably stop without missing relevant records or fail to
utilize the full potential work-savings. Only one method reliably meets
the set recall target, but stops conservatively.}
An independent test of the stop criteria offered with one widely used
open tool, in 35,000 simulated screenings of five health datasets,
reports that \textcquote[abstract]{Kempny2026}{\textins*{n}one of the
tested stop criteria could consistently identify all relevant studies
across all datasets}; the share of a dataset that had to be screened to
find all relevant studies ranged from 2.9\,\% to 76.9\,\% on average.
A heuristic that stops after 50 irrelevant records in a row missed its
target of 95\,\% recall in 39\,\% of the cases in the simulations of
\textcite[p.~11]{Callaghan2020}.
A live mental-health review, reported in a preprint read only in its
abstract, concludes that \textcquote[abstract]{Slade2026}{stopping rules
risk missing relevant studies}.
For a ranked web source, a study read only in its abstract found that
\textcquote[abstract]{Briscoe2023}{the lowest ranking relevant studies
were ranked 227th and 215th for the two systematic reviews}, so the first
100 results of Google Search would have missed them.

\paragraph{A statistical criterion holds its target, at a cost.}
The one kind of rule that held a stated recall target in these
comparisons tests, on a random sample of the records not yet seen,
whether the target has been reached:
\textcquote[abstract]{Callaghan2020}{The stopping criteria are shown on
test datasets to achieve a reliable level of recall, while still
providing work reductions of on average 17\,\%. Other methods proposed
previously are shown to provide inconsistent recall and work reductions
across datasets.}
In ten health reviews, stopping a prioritised screen at the point where
95\,\% of the relevant records had in fact been found cut the screening
by a median of 47.1\,\%, and \textcquote[abstract]{Hamel2020}{\textins*{a}mong
the 100 simulations (10 simulations per review), none of these 5\,\% of
records were a final included study in the systematic review}; but that
stopping point was known only afterwards, not estimated during the
screen.
Using a tool's relevance predictions in place of one of two reviewers
missed a median of 1\,\% (0 to 14\,\%) of relevant records in eleven
systematic reviews, and in six rapid reviews
\textcquote[abstract]{Gates2020}{\textins*{t}he median (range)
proportion missed records for both approaches was 6 (0--22)\,\%}.

\paragraph{How well a rule does depends on the dataset.}
In 27 reviews in education and educational psychology, heuristic rules
retrieved 95\,\% of the relevant abstracts, but
\textcquote[abstract]{Campos2024}{the performance of the heuristic
stopping criteria depended on the learning algorithm used and the length
and proportion of relevant papers in an abstract collection}.
An evaluation in psychology, read only in its abstract, found
\textcquote[abstract]{Koenig2024}{performance differences between
stopping rules regarding all performance measures and variations in the
performance of stopping rules across different prevalence ratios}.
The authors of a conservative heuristic that combines several rules leave
the decision to the reviewer:
\textcquote[abstract]{Boetje2024}{the decision to stop the active
learning process depends on careful consideration of the trade-off
between the costs of additional record labeling against the potential
errors of the current model for the specific dataset and context}.
A commentary by the authors of the statistical criterion, with
methodologists of evidence synthesis, states the general point:
\textcquote[p.~2]{Callaghan2024}{\textins*{w}ithout prior knowledge of
the number of relevant records, any approach to stopping early risks
missing relevant studies that would have been retrieved had all records
been screened}.

\section{What a sample of the rest can say}
\label{sec:stopping-sample}

The statistical criterion rests on
\textcquote[p.~5]{Callaghan2020}{drawing a random sample from the
remaining unseen documents}.
The procedure of \cref{ch:procedure} uses the same kind of evidence for a
smaller purpose.
It does not stop when a recall target is reached; it screens a tier
chosen in advance, draws a seeded random sample of what the tier left
out, and reports the share of bearing records in the sample with its
interval.
That share estimates how much of the unscreened rest bears on the
question; it does not say which records those are, and a sample in which
none bears does not show that none does.
No source found tests such a sample against a full screen; the reason to
draw it is that it is the only statement about the unscreened records
that the sources above leave room for.

\section{Discussion and limitations}
\label{sec:stopping-discussion}

The evidence is mostly simulation on collections of reviews already
screened in full, most of them in health research; one evaluation is in
education, one in psychology, and one live review is reported only in a
preprint's abstract.
The one result in which no included study was lost at the stopping point
used a point known only in retrospect.
None of the sources measures whether the studies missed at a stopping
point would have changed a review's conclusions.
The psychology evaluation and the study of Google Search could be read
only in their abstracts (the publisher refused the full texts), and the
capture-recapture stopping rule, one of the named studies, was not
obtained and is not cited.
The search itself had one screener, an agent (\cref{sec:search-limits}).

\section{Conclusion}
\label{sec:stopping-conclusion}

No: in the comparisons found, most stopping rules either missed relevant
records or saved little work, and only a statistical criterion built on
a random sample of the unseen records held its recall target, saving
about a sixth of the screening on test collections; how well any rule
did depended on the dataset.
For the procedure this means that a tier is not a stopping rule.
It decides what is read first; the records it leaves out are counted and
sampled, and the method says so (\cref{ch:procedure}).

\section{Future work}
\label{sec:stopping-future}

This search left undone the full texts of the psychology evaluation and
the study of Google Search, and the capture-recapture rule, which a
reader with access to the publishers can add.
The field lacks a test of stopping rules outside health and education,
in computing in particular, and a measurement of whether the studies a
rule misses change what a review concludes; either would say how much a
tier can be trusted to stand in for a full screen.

\section{The records judged bearing}
\label{sec:stopping-hits}

% Author's note: table by generate_audit_table on a copy of screenlit-q1,
% bearing rows only, cited = cited in doc/bibliography.bib.
% --- scholar LaTeX fragment --------------------------------------------
% Generated by: scholar sessions export -f table
% Session: screenlit-q1
% \input this file from the body of a document whose
% preamble loads:
%   \usepackage{longtable}
%   \usepackage{booktabs}
%   \usepackage{array}
% ----------------------------------------------------------------------

\begingroup\footnotesize
\begin{longtable}{@{}%
  >{\raggedright\arraybackslash}p{0.44\textwidth}c%
  >{\raggedright\arraybackslash}p{0.13\textwidth}%
  >{\raggedright\arraybackslash}p{0.26\textwidth}@{}}
\caption{Records bearing on the claim, Question 1, stopping a ranked screen early (11 cited, 22 supporting and 20 qualifying or contradicting it, out of 897 unique hits; the 165 adjacent hits, the 282 false hits and 397 unclassified are not listed). The last column gives the reason for each inclusion; for rows decided by an agent the agent's confidence is given. A row is cited when the text cites the work, whoever decided on its records. Databases searched: OpenAlex (OA; not open access), IEEE Xplore (IEEE), Web of Science (WoS), Scopus, Semantic Scholar (S2), arXiv. Searched 2026-11-23.}\label{tab:screenlit-q1}\\
\toprule
Title & Year & Database & Reason \\
\midrule
\endfirsthead
\multicolumn{4}{@{}l}{\emph{\tablename~\thetable{} (cont.)}}\\
\toprule
Title & Year & Database & Reason \\
\midrule
\endhead
\midrule \multicolumn{4}{r@{}}{\emph{continued on next page}}\\
\endfoot
\bottomrule
\endlastfoot
An evaluation of DistillerSR's machine learning-based prioritization tool for title/abstract screening -- impact on reviewer-relevant outcomes & 2020 & OA, WoS, Scopus, S2 & \emph{cited}: supports the claim (0.70) \\
An evaluation of the performance of stopping rules in AI-aided screening for psychological meta-analytical research & 2024 & Scopus, WoS & \emph{cited}: qualifies or contradicts the claim (0.70) \\
Computer-assisted screening in systematic evidence synthesis requires robust and well-evaluated stopping criteria & 2024 & OA, Scopus & \emph{cited}: qualifies or contradicts the claim (0.60) \\
Don't Stop Me Now, `Cause I'm Having a Good Time Screening: Evaluation of Stopping Methods for Safe Use of Priority Screening in Systematic Reviews & 2026 & OA, S2 & \emph{cited}: qualifies or contradicts the claim (0.75) \\
Feasibility and desirability of screening search results from Google Search exhaustively for systematic reviews: A crosscase analysis & 2023 & OA, WoS, Scopus, S2 & \emph{cited}: qualifies or contradicts the claim (0.60) \\
Machine learning-assisted versus conventional title and abstract screening for a mental health systematic review (Preprint) & 2026 & OA & \emph{cited}: qualifies or contradicts the claim (0.70) \\
Screening Smarter, Not Harder: A Comparative Analysis of Machine Learning Screening Algorithms and Heuristic Stopping Criteria for Systematic Reviews in Educational Research & 2024 & OA, Scopus, S2, WoS & \emph{cited}: qualifies or contradicts the claim (0.70) \\
Statistical stopping criteria for automated screening in systematic reviews & 2020 & OA, WoS, Scopus, S2 & \emph{cited}: qualifies or contradicts the claim (0.80) \\
The SAFE procedure: a practical stopping heuristic for active learning-based screening in systematic reviews and meta-analyses & 2024 & WoS, Scopus, S2, OA & \emph{cited}: qualifies or contradicts the claim (0.60) \\
The semi-automation of title and abstract screening: a retrospective exploration of ways to leverage Abstrackr's relevance predictions in systematic and rapid reviews & 2020 & OA, S2, WoS, Scopus & \emph{cited}: qualifies or contradicts the claim (0.60) \\
When to stop reviewing: validation of stop criteria in ASReview & 2026 & OA, Scopus, WoS, S2 & \emph{cited}: qualifies or contradicts the claim (0.70) \\
A text-mining tool generated title-abstract screening workload savings: performance evaluation versus single-human screening & 2022 & OA, S2, WoS, Scopus & supports the claim (0.70) \\
AI-supported screening for systematic literature reviews informing EULAR recommendations and points to consider & 2026 & OA, S2, Scopus & supports the claim (0.75) \\
Can artificial intelligence separate the wheat from the chaff in systematic reviews of health economic articles? & 2023 & OA, WoS, Scopus, S2 & supports the claim (0.55) \\
Confidence-Based Stopping Methods for Systematic Reviews & 2026 & arXiv, WoS, OA, S2, Scopus & supports the claim (0.65) \\
Decision-Theoretic Stopping Rules for Document Screening & 2026 & OA, arXiv, WoS, Scopus, S2 & supports the claim (0.55) \\
Enhancing recall in automated record screening: A resampling algorithm & 2024 & OA, WoS, Scopus, S2 & supports the claim (0.50) \\
Improving Stopping Methods for Technology Assisted Review & 2026 & OA & supports the claim (0.55) \\
Increasing the Efficiency of High-Recall Information Retrieval & 2019 & OA, S2 & supports the claim (0.55) \\
Increasing the efficiency of study selection for systematic reviews using prioritization tools and a single-screening approach & 2023 & OA, S2, WoS, Scopus & supports the claim (0.60) \\
Machine learning enables automated screening for systematic reviews and meta-analysis in urology & 2024 & OA, Scopus, S2 & supports the claim (0.50) \\
Machine Learning Reduced Workload for COVID-19 Literature & 2025 & OA & supports the claim (0.50) \\
Natural language processing was effective in assisting rapid title and abstract screening when updating systematic reviews & 2021 & WoS, Scopus & supports the claim (0.60) \\
Research Screener: a machine learning tool to semi-automate abstract screening for systematic reviews & 2021 & OA, S2, WoS, Scopus & supports the claim (0.60) \\
RLStop: A Reinforcement Learning Stopping Method for TAR & 2024 & arXiv, WoS & supports the claim (0.60) \\
Sensitivity and specificity of alternative screening methods for systematic reviews using text mining tools & 2023 & WoS, Scopus & supports the claim (0.55) \\
Stopping Methods for Technology-assisted Reviews Based on Point Processes & 2023 & OA, arXiv, WoS, Scopus, S2 & supports the claim (0.65) \\
SWIFT-Active Screener: Accelerated document screening through active learning and integrated recall estimation & 2020 & OA, WoS, Scopus, S2 & supports the claim (0.65) \\
Text mining to support abstract screening for knowledge syntheses: a semi-automated workflow & 2021 & OA, S2, WoS, Scopus & supports the claim (0.55) \\
The capture--mark--recapture technique can be used as a stopping rule when searching in systematic reviews & 2008 & OA, WoS, Scopus, S2 & supports the claim (0.45) \\
Usefulness of machine learning softwares to screen titles of systematic reviews: a methodological study & 2023 & OA, S2, WoS, Scopus & supports the claim (0.50) \\
Using Chao's Estimator as a Stopping Criterion for Technology-Assisted Review & 2024 & arXiv, OA, WoS, Scopus, S2 & supports the claim (0.60) \\
Validity of using a semi-automated screening tool in a systematic review assessing non-specific effects of respiratory vaccines & 2025 & OA, S2, WoS, Scopus & supports the claim (0.50) \\
A Study on a Stopping Strategy for Systematic Reviews Based on a Distributed Effort Approach & 2020 & OA, WoS, Scopus & qualifies or contradicts the claim (0.45) \\
Application of Machine Learning Classifiers in Rapid Reviews for Health Research: A Case Example Using EPPIReviewer & 2026 & OA & qualifies or contradicts the claim (0.45) \\
Assessing the Accuracy of Machine-Assisted Abstract Screening With DistillerAI: A User Study & 2019 & OA & qualifies or contradicts the claim (0.65) \\
Assessing the accuracy of machine-assisted abstract screening with DistillerAI: a user study & 2019 & OA, S2, WoS, Scopus & qualifies or contradicts the claim (0.60) \\
Can machine learning help accelerate article screening for systematic reviews? Yes, when article separability in embedding space is high & 2025 & Scopus & qualifies or contradicts the claim (0.65) \\
Decoding semi-automated title-abstract screening: findings from a convenience sample of reviews & 2020 & OA, S2, WoS, Scopus & qualifies or contradicts the claim (0.55) \\
Faster title and abstract screening? Evaluating Abstrackr, a semi-automated online screening program for systematic reviewers & 2015 & OA, S2, WoS, Scopus & qualifies or contradicts the claim (0.55) \\
Machine learning for screening prioritization in systematic reviews: Comparative performance of Abstrackr and EPPI-Reviewer & 2020 & Scopus & qualifies or contradicts the claim (0.50) \\
Measuring the impact of screening automation on meta-analyses of diagnostic test accuracy & 2019 & OA, WoS, Scopus, S2 & qualifies or contradicts the claim (0.55) \\
Optimising screening efficiency in evidence synthesis on health Technology: A simulation study using ASReview & 2026 & WoS, Scopus, OA, S2 & qualifies or contradicts the claim (0.60) \\
Performance and Usability of Machine Learning for Screening in Systematic Reviews: A Comparative Evaluation of Three Tools & 2019 & OA, S2, WoS, Scopus & qualifies or contradicts the claim (0.55) \\
Performance of Two AI Approaches in ASReview Compared With Manual Screening for Dementia Care Literature Screening: Comparative Analysis & 2026 & OA, WoS, Scopus, S2 & qualifies or contradicts the claim (0.70) \\
Prioritising references for systematic reviews with RobotAnalyst: A user study & 2018 & OA, Scopus, S2, WoS & qualifies or contradicts the claim (0.55) \\
SnowBallSLR: saturation-aware and reproducible citation searching for systematic reviews & 2026 & OA, Scopus & qualifies or contradicts the claim (0.45) \\
Strategizing AI utilization for psychological literature screening: A comparative analysis of machine learning algorithms and key factors to consider & 2026 & Scopus, OA, S2, WoS & qualifies or contradicts the claim (0.50) \\
Technology-assisted title and abstract screening for systematic reviews: a retrospective evaluation of the Abstrackr machine learning tool & 2018 & OA, S2, WoS, Scopus & qualifies or contradicts the claim (0.50) \\
The efficiency of machine learning-assisted platform for article screening in systematic reviews in orthopaedics & 2022 & OA, WoS, Scopus, S2 & qualifies or contradicts the claim (0.50) \\
To include or not to include? A prescription from the pharmacy on how to use active learning-assisted screening in systematic reviews & 2026 & OA, WoS, Scopus, S2 & qualifies or contradicts the claim (0.45) \\
Unsupervised title and abstract screening for systematic review: a retrospective case-study using topic modelling methodology & 2023 & OA, S2, WoS, Scopus & qualifies or contradicts the claim (0.50) \\
Write your abstracts carefully---The impact of abstract reporting quality on findability by semi-automated title-abstract screening tools & 2025 & OA, S2, WoS, Scopus & qualifies or contradicts the claim (0.55) \\
\end{longtable}
\endgroup

